\begin{table}[htbp]
\centering
\small
\setlength{\tabcolsep}{4pt}
\begin{tabular}{@{}llccl@{}}
\toprule
\multicolumn{2}{c}{结构赋值类型}&函数条件&DAG可识别&参见\\
\midrule
（一般）SCM&$X_j:=f_j(X_{\PA_j},N_j)$&—&否&命题7.1\\
ANM&$X_j:=f_j(X_{\PA_j})+N_j$&非线性&是&定理7.7（i）\\
CAM&$X_j:=\sum_{k\in\PA_j}f_{jk}(X_k)+N_j$&非线性&是&定理7.7（ii）\\
线性高斯&$X_j:=\sum_{k\in\PA_j}\beta_{jk}X_k+N_j$&线性&否&习题7.13\\
线性高斯、误差等方差&$X_j:=\sum_{k\in\PA_j}\beta_{jk}X_k+N_j$&线性&是&命题7.5\\
\bottomrule
\end{tabular}
\caption{高斯噪声下若干已知可识别性结果的汇总。对于非高斯噪声，也有可识别性结果，但其技术细节更复杂。}
\label{tab:ch07-identifiability}
\end{table}
